%==================================================================
% BAB 2 — TINJAUAN PUSTAKA
%==================================================================

\chapter{TINJAUAN PUSTAKA}

\section{Tinjauan Pustaka}

% CATATAN PENTING
% Gunakan MATRIKS SINTESIS, bukan tabel per-penelitian dengan kolom
% "perbedaan dengan penelitian ini". Pada matriks, celah penelitian
% TERLIHAT sebagai pola (kolom yang kosong di semua baris), sehingga
% Anda cukup menuliskan apa yang tampak — tidak perlu mengarang kebaruan.
%
% Kolom dipilih berdasarkan relevansinya dengan rumusan masalah Anda.
% Empat sampai enam penelitian yang benar-benar dibedah jauh lebih baik
% daripada lima belas yang sekadar didaftar.

Bagian ini menyintesiskan penelitian terdahulu yang relevan dengan metode dan 
teknik yang digunakan dalam perancangan Sistem Informasi Manajemen Magang (SIMA) 
pada Program Studi Sistem Informasi Unika Atma Jaya. Perbandingan disajikan pada 
\cref{tab:sintesis} dengan kolom pembanding yang dipilih berdasarkan keterkaitannya dengan 
rumusan masalah penelitian ini, yaitu metode pengembangan, mekanisme kontrol akses, 
serta teknik pengujian kelayakan sistem.

\small
\begin{xltabular}{\textwidth}{
  >{\RaggedRight\hsize=0.85\hsize}X 
  >{\RaggedRight\hsize=0.95\hsize}X 
  >{\RaggedRight\hsize=0.90\hsize}X 
  >{\RaggedRight\hsize=0.90\hsize}X 
  >{\RaggedRight\hsize=1.40\hsize}X
}
% --- Judul & Keterangan Tabel ---
\caption{Matriks sintesis penelitian terdahulu} \label{tab:sintesis} \\
\toprule
\textbf{Penelitian} & \textbf{Konteks pengguna} & \textbf{Instrumen} & \textbf{Teknik analisis} & \textbf{Temuan utama} \\
\midrule
\endfirsthead

% --- Judul di Halaman Sambungan (Halaman Berikutnya) ---
\multicolumn{5}{c}{{\bfseries Tabel 2.1} -- Lanjutan dari halaman sebelumnya} \\[0.5em]
\toprule
\textbf{Penelitian} & \textbf{Konteks pengguna} & \textbf{Instrumen} & \textbf{Teknik analisis} & \textbf{Temuan utama} \\
\midrule
\endhead

% --- Garis Penutup Halaman Pertama ---
\midrule
\endfoot

% --- Garis Penutup Halaman Terakhir ---
\bottomrule
\endlastfoot

% --- Isi Data Tabel ---
\cite{Parlika2023} & 16 studi penerapan model RAD (2014--2022) pada berbagai skala proyek & Narrative Literature Review & Deskriptif kualitatif & Model RAD menghemat waktu pengembangan 21--120 hari dibanding model tradisional; cocok untuk sistem yang dapat dipecah menjadi modul \\
\midrule
\cite{Mintarsih2023} & Sistem Informasi Perpustakaan berbasis web (SMC Foundation) & Black-box testing teknik state transition & Pengujian transisi alur antar-menu & 6 dari 8 transisi berhasil sesuai harapan; 2 transisi gagal karena pembatasan hak akses; ditemukan bug validasi tipe data pada form edit \\
\midrule
\cite{AliyahAliyah2025} & Sistem informasi pengelolaan keuangan dan inventaris (CV. Nitah Tirta Makassar) & User Acceptance Testing, kuesioner skala Likert 1--5 & Deskriptif, perhitungan skor interpretasi & Rata-rata skor UAT 87 persen (kategori sangat baik) pada 4 variabel: fungsionalitas, kinerja, antarmuka, efisiensi \\
\midrule
\cite{Fiandhika2024} & Sistem e-budgeting pada BUMN sektor asuransi & System Usability Scale + Heuristic Evaluation + wawancara & Studi komparatif skor antarmetode & Skor SUS tidak sepenuhnya valid bila dibandingkan evaluasi ahli; hanya indikator ease of use yang konsisten valid di seluruh metode \\
\midrule
\cite{SafiraArmah2024} & Konsep umum Sistem Informasi Manajemen pada organisasi/perusahaan & Studi literatur & Deskriptif & SIM merupakan perpaduan manusia, teknologi, dan prosedur; manfaat utama pada aksesibilitas data dan efisiensi pengambilan keputusan \\
\midrule
\cite{MSahyudi2025} & Implementasi RBAC pada aplikasi Enterprise Resource Planning & Systematic Literature Review (protokol PRISMA) & Sintesis tematik dan komparatif & Integrasi RBAC dengan AI meningkatkan keamanan hingga 37 persen; otomatisasi RBAC menurunkan biaya operasional 42 persen dan waktu manajemen akses 65 persen \\
\end{xltabular}

% Paragraf berikut MEMBACA POLA pada matriks, lalu menyatakan posisi
% penelitian. Perhatikan bahwa celah dinyatakan berdasarkan apa yang
% terlihat pada tabel, bukan diklaim begitu saja.
Berdasarkan \cref{tab:sintesis} terlihat dua pola. Pertama, penelitian-penelitian 
terdahulu yang membahas model pengembangan (RAD), kontrol akses (RBAC), dan teknik 
pengujian (black-box, UAT, SUS) umumnya dilakukan secara terpisah pada satu aspek 
sistem saja, dan belum ada yang mengombinasikan kelima pendekatan tersebut dalam 
satu alur pengembangan sistem informasi yang utuh. Kolom objek/konteks pada matriks 
memperlihatkan bahwa objek penelitian terdahulu didominasi oleh sistem informasi 
bisnis dan organisasi (perpustakaan, keuangan, e-budgeting, ERP), sedangkan konteks 
sistem informasi manajemen magang pada perguruan tinggi dengan multi-aktor 
(mahasiswa, Kaprodi, Koordinator Magang, Dosen Pembimbing, Pembimbing Eksternal) 
belum banyak dieksplorasi. Kedua, temuan \cite{Fiandhika2024} 
mengindikasikan bahwa pengukuran usability dengan SUS saja berpotensi tidak 
sepenuhnya valid tanpa didukung metode evaluasi lain, sehingga penggunaan SUS 
sebaiknya dilengkapi pengujian fungsional seperti black-box testing dan UAT agar 
hasil evaluasi kelayakan sistem lebih komprehensif.

Perbedaan pola penggunaan metode secara terpisah tersebut berpotensi menyebabkan 
evaluasi kelayakan sistem menjadi tidak menyeluruh, sebagaimana dikemukakan oleh 
\cite{Fiandhika2024}. Penelitian ini karena itu menguji apakah kombinasi 
model RAD untuk pengembangan, RBAC untuk kontrol akses, serta black-box testing, 
UAT, dan SUS untuk pengujian kelayakan, dapat diterapkan secara terintegrasi pada 
konteks sistem informasi manajemen magang di lingkungan perguruan tinggi, sekaligus 
melengkapi evaluasi usability dengan pengujian fungsional agar hasil yang diperoleh 
lebih valid.

\section{Landasan Teori}
% Susun subbab sesuai konsep utama penelitian Anda. Tiap subbab memuat
% definisi, unsur pokok, dan rujukannya. Hindari menyalin definisi
% panjang; ringkas dengan bahasa sendiri lalu sitasi sumbernya.

\subsection{Sistem Informasi Manajemen}
Sistem Informasi Manajemen (SIM) adalah sistem yang dirancang untuk menyediakan 
informasi bagi para manajer dalam rangka mendukung operasi, manajemen, dan 
pengambilan keputusan dalam sebuah organisasi, dengan mengintegrasikan berbagai 
sumber data menjadi informasi yang bermakna. SIM merupakan perpaduan dari tiga 
komponen utama, yaitu manusia, teknologi, dan prosedur (\cite{SafiraArmah2024}).

\subsection{Rapid Application Development (RAD)}
Rapid Application Development (RAD) merupakan metodologi pengembangan perangkat 
lunak yang hanya memerlukan perencanaan minimal untuk mempercepat proses 
pengembangan melalui tiga tahap, yaitu Requirement Planning, RAD Design Workshop, 
dan Implementation. Model RAD lebih menitikberatkan kegiatan tahap pengembangan 
daripada tahap perencanaan dan cocok digunakan untuk proyek yang dapat 
dipisah-pisahkan menjadi beberapa modul (\cite{Parlika2023}).

Tahapan Pelaksanaan Model RAD
\begin{daftarnomor}
    \item Requirement Planning \\
          Pada fase ini, system analyst, customer, dan user akan bertemu dan bersama-sama merumuskan tujuan serta requirements 
          dari perangkat lunak yang akan dikembangkan. Waktu pada tahapan perencanaan ini hanya dialokasikan secara singkat, 
          karena proses pengembangan aplikasi akan berjalan secara fleksibel sesuai kebutuhan, berbeda dengan model tradisional 
          yang mengalokasikan waktu perencanaan lebih panjang.
    \item RAD Design Workshop \\
          Pada fase ini akan dibuat sketsa dan diagram yang nantinya akan dijadikan panduan dalam pengembangan perangkat lunak. 
          ERD, UML, dan diagram-diagram lain dirancang oleh system analyst bersama dengan programmer pada tahap ini. Di sini juga 
          dibuat rancangan tampilan antarmuka dari perangkat lunak yang dibuat. Fase ini merupakan tahap kolaboratif di mana 
          desain sistem disusun secara langsung bersama pengguna untuk memastikan kesesuaian dengan kebutuhan mereka.
    \item Implementation \\
          Kode dari perangkat lunak akan dibuat pada fase ini. Akan pula dilakukan pengujian oleh QA dan oleh user untuk 
          memutuskan apakah sistem sudah layak untuk digunakan. Fase ini adalah tahapan memperkenalkan sistem kepada user interface. 
          Dalam fase ini, terutama dalam tahap pengujian, user akan semakin sering dilibatkan dalam proses pengembangan perangkat lunak. 
          Tahap ini bisa dikatakan sebagai tahapan terakhir dalam model RAD, karena jika perangkat lunak telah dinilai layak oleh 
          user dan QA dari proses pengujian, maka perangkat lunak dapat dirilis dan digunakan.
\end{daftarnomor}

\subsection{Role-Based Access Control (RBAC)}
Role-Based Access Control (RBAC) adalah model keamanan yang membatasi akses sistem 
kepada pengguna berdasarkan peran (role) mereka dalam organisasi, terdiri atas tiga 
komponen utama yaitu pengguna (users), peran (roles), dan izin (permissions), 
dengan menerapkan prinsip least privilege di mana pengguna hanya diberikan akses 
minimal yang diperlukan untuk menjalankan tugasnya (\cite{MSahyudi2025}).

\subsection{Black-Box Testing}
Black-box testing adalah pengujian fungsional yang memverifikasi hasil eksekusi 
aplikasi berdasarkan masukan yang diberikan untuk memastikan fungsionalitas 
sistem sesuai dengan persyaratan, tanpa menguji berdasarkan source code program. 
Salah satu teknik black-box testing adalah state transition testing, yang 
dilakukan dengan melihat kesesuaian alur dari satu state ke state berikutnya 
sesuai transisi, keadaan, dan kejadian di antara inputan (\cite{Mintarsih2023}).

\subsection{User Acceptance Testing (UAT)}
User Acceptance Testing (UAT) merupakan salah satu tahap penting dalam pengujian 
perangkat lunak yang dilakukan oleh pengguna akhir (end-user) untuk memastikan 
bahwa perangkat lunak yang dikembangkan memenuhi kebutuhan dan harapan pengguna 
sebelum diluncurkan secara resmi. Pengujian UAT umumnya difokuskan pada beberapa 
variabel utama, yaitu fungsionalitas sistem, kinerja sistem, pengalaman dan 
tampilan antarmuka, serta efisiensi dan produktivitas (\cite{AliyahAliyah2025}).

\subsection{System Usability Scale (SUS)}
System Usability Scale (SUS) adalah metode yang banyak digunakan untuk 
mengevaluasi usability suatu sistem secara cepat melalui 10 pernyataan 
dengan skala Likert. Meskipun demikian, skor SUS lebih merepresentasikan 
persepsi pengguna dan tidak selalu mencerminkan usability aktual secara 
menyeluruh, sehingga disarankan untuk dilengkapi dengan metode evaluasi 
lain seperti heuristic evaluation agar hasil pengukuran lebih dapat 
diverifikasi (\cite{Fiandhika2024}).

\section{Justifikasi Pemilihan Metode}
% SUBBAB INI SERING DILUPAKAN, PADAHAL PALING PENTING.
% Deskriptor KKNI jenjang 6 menuntut kemampuan memilih alternatif solusi
% secara sadar. Di sinilah kemampuan itu dibuktikan.
% Bandingkan minimal dua alternatif, lalu nyatakan alasan pemilihan Anda.

Terdapat beberapa model pengembangan perangkat lunak yang lazim digunakan, 
di antaranya Waterfall dan Prototyping. Model Waterfall bersifat linear sehingga 
pengelolaannya mudah karena seluruh kebutuhan telah diidentifikasi dan 
didokumentasikan di awal, namun tidak fleksibel terhadap perubahan kebutuhan 
dan membutuhkan waktu pengembangan yang relatif lama hingga sistem dapat 
digunakan oleh pengguna. Model Prototyping mengikuti pendekatan evolutionary 
dan cocok digunakan ketika kebutuhan pengguna belum diputuskan, tetapi menambah 
kompleksitas sistem seiring bertambahnya waktu pengembangan serta menimbulkan 
biaya tambahan untuk pembuatan prototipe awal (\cite{Parlika2023}).

RAD dipilih sebagai model pengembangan dalam penelitian ini karena sistem 
Manajemen Magang dapat dipecah menjadi beberapa modul fungsional 
(pendaftaran, bimbingan, penilaian) yang relatif independen, serta melibatkan 
pengguna secara langsung dalam setiap tahap pengembangan sehingga produk 
yang dihasilkan lebih sesuai dengan kebutuhan Program Studi Sistem Informasi 
Unika Atma Jaya. Selain itu, RAD terbukti dapat menghemat waktu pengembangan 
dibandingkan model tradisional, yang sesuai dengan kebutuhan penelitian ini 
yang memiliki keterbatasan waktu penyusunan tugas akhir (\cite{Parlika2023}).

Untuk kontrol akses, pendekatan manual tanpa model baku berisiko menimbulkan 
ketidakkonsistenan hak akses antar-aktor, mengingat sistem ini melibatkan 
lima aktor dengan kewenangan berbeda. RBAC dipilih karena mampu membatasi 
akses setiap aktor sesuai perannya masing-masing berdasarkan prinsip least 
privilege, sehingga mahasiswa, Kaprodi, Koordinator Magang, Dosen Pembimbing, 
dan Pembimbing Eksternal hanya dapat mengakses fitur yang relevan dengan tugas 
mereka (\cite{MSahyudi2025}).

Untuk pengujian kelayakan sistem, alternatif yang tersedia adalah menggunakan 
satu metode evaluasi saja, baik black-box testing, UAT, maupun SUS. Namun, 
temuan \cite{Fiandhika2024} menunjukkan bahwa skor SUS semata tidak 
selalu valid tanpa dukungan metode lain. Oleh karena itu, penelitian ini 
menggunakan kombinasi tiga metode: black-box testing dengan teknik state 
transition untuk memvalidasi kesesuaian alur fungsional antarmenu 
(\cite{Mintarsih2023}), UAT untuk mengukur penerimaan pengguna akhir terhadap 
fungsionalitas dan kinerja sistem (\cite{AliyahAliyah2025}), dan SUS 
untuk mengukur tingkat usability secara umum. Kombinasi ketiga metode ini 
dipilih agar evaluasi kelayakan sistem mencakup baik aspek fungsional 
(alur kerja sistem) maupun aspek penerimaan dan pengalaman pengguna secara 
lebih menyeluruh.
